% Entropie des échelles E8M0 des 72 tenseurs d'experts de gpt-oss-20b, selon le contexte :
% ordre 0, sachant l'échelle du bloc de gauche, sachant la médiane de la ligne (coût de la
% médiane compris). Données : fondations/journaux/03-entropie/csv/fig-echelles-gptoss20b.csv
% (vers_csv.py, depuis mxfp4-contexte-gptoss20b.json).
% Compilation seule : lualatex fig-echelles.tex (depuis ce dossier).
\documentclass[tikz,border=4pt]{standalone}
\usepackage{fontspec}
\usepackage{pgfplots}
\usepackage{pgfplotstable}
\pgfplotsset{compat=1.18}
\providecommand{\donnees}{../../../fondations/journaux/03-entropie/csv}
\begin{document}
\pgfplotstableread[col sep=comma]{\donnees/fig-echelles-gptoss20b.csv}\table
\begin{tikzpicture}
\begin{axis}[
  width=13cm, height=6.5cm, xmin=-1, xmax=72, ymin=0, ymax=2.5,
  xlabel={tenseur d'experts (couche × \{down, gate, up\})}, ylabel={bits par échelle (sur 8)},
  legend style={at={(0.5,-0.22)}, anchor=north, font=\small, legend columns=3},
  ymajorgrids, grid style={gray!25}, mark size=1.2pt,
]
\addplot[only marks, mark=o, gray] table[x=indice, y=H_e] {\table};
\addplot[only marks, mark=triangle*, orange!80!black] table[x=indice, y=H_e_gauche] {\table};
\addplot[only marks, mark=*, blue!70!black] table[x=indice, y=H_e_mediane] {\table};
\legend{ordre 0, sachant le bloc de gauche, sachant la médiane de ligne}
\end{axis}
\end{tikzpicture}
\end{document}
